\chapter{Introduction}
\label{chap:introduzione}

\section{Background}

A backup, in the field of cybersecurity, is a disaster recovery process: a procedure
that makes it possible to recover lost data, for example following an attack, a hardware or
software failure, or a human error. One of the first technical responses to data loss is the redundancy of the disks
on which the data is stored: in 1988 Patterson, Gibson and Katz formalized RAID (Redundant Array of
Inexpensive Disks) \cite{patterson-gibson-katz-raid-1988}, which distributes the same data across multiple
inexpensive disks to obtain a reliability comparable to that of a single, more expensive disk.

This redundancy, however, remains local and does not protect against an event that destroys the entire installation. From
this arises the need for a remote backup: a copy of the data kept in a place other than that
of the originals. The simplest way is to entrust it to third parties, an acquaintance who owns the hardware or
a business that offers this kind of service, but this requires trusting that the data
will not be read, modified, or sold by whoever keeps it.

Whoever offers the service, for their part, finds it convenient to serve multiple users at the same time. This,
however, requires distinguishing between users, and guaranteeing that each one accesses only their own data: in
practice, it requires an authentication and access control system.

Among the various existing authentication systems, the one adopted and discussed in this thesis is based on
email and password which, being sensitive data, require a high level of protection.
The principle on which this kind of authentication rests is simple:
the user registers by sending their credentials to the server; during login they send them again, and the
server compares them with those saved at registration. If they match, the user is authenticated.

The fact that the server holds the credentials, however, exposes all users to the risk of their being
stolen. Numerous systems exist to prevent this; one of these, examined in depth in this thesis, is the use of
a hash function.
The server, instead of saving the credentials in plaintext, saves them
in the form of a hash. During login, it recomputes the hashes of the credentials provided by the user and
checks whether they match the saved ones.

\subsection{Hash functions}

Hash functions are non-invertible functions that transform an input of arbitrary
length into an output of fixed length, strictly correlated with the input data, called a hash.
Denoting a hash function by $H$, it enjoys the following fundamental properties:

\begin{itemize}
    \item \textbf{preimage resistance:} given $h$, it is computationally intractable to find $x$
    such that $H(x) = h$;
    \item \textbf{second-preimage resistance:} given $x$, it is computationally intractable to
    find $y \neq x$ such that $H(x) = H(y)$;
    \item \textbf{collision resistance:} it is computationally intractable to find any pair
    $(x, y)$ with $x \neq y$ such that $H(x) = H(y)$.
\end{itemize}


\section{Problem statement}

Among the various causes of data loss, the project focuses on the one arising from an attack by a
malicious actor, adopting the principles of zero-trust and defense-in-depth to guarantee the
security of the data even if any component of the system is compromised.

\section{Objectives}

RAM-USB, Remotely Accessible Multi-User Secure Backup, is an in-depth case study on the design,
implementation and deployment of a backup service that satisfies the requirements
described in \autoref{chap:requisiti}. It is not meant to compete with existing commercial
solutions, but to discuss its own design choices, illustrating their pros and cons
transparently, with a security-first approach focused on defending the system against external attacks.


\section{Methodology}
Some of the components of RAM-USB originate from the brainstorming phase carried out together with Riccardo
Gottardi \cite{gottardi-github}.

The adopted methodology combines the spiral model for the definition and evolution of the
requirements \cite{Boehm88} with the Test-Driven Development approach for the implementation and the V-model for
their verification \cite{ramusb-contributing}.

The language used for the entire stack is Go, a compiled language explicitly designed to handle
many simultaneous network connections efficiently.

Bash was used to orchestrate the infrastructure around the Go code (starting the services, supervising the processes in the
containers, sequencing at startup).

The code is versioned with Git and published free of charge, in its entirety, on GitHub, distributed under
the MIT license.

The implementation was carried out entirely with the assistance of the Large Language Models Claude
Sonnet 4.5 and Sonnet 5, in multi-agent form through the Claude Code harness running in a virtual
machine on Proxmox, for security and privacy reasons.

Generative artificial intelligence was employed exclusively as a tool under close
human supervision.
The design, the structure of the code, the module boundaries, the signatures, the behavior of the individual
functions, and the corresponding tests remain a human decision; the assistant implemented only what was
specified, in small, incremental, and verifiable steps, exercising latitude only over low-level
implementation details.
Every change was reviewed and explicitly approved before being introduced into the codebase.
The numerous strategies for mitigating implementation errors and hallucinations of the two models, as well as the
management of the multiple agents used, are not the subject of this thesis.

\section{Thesis structure}

\autoref{chap:principi-fondanti} examines in depth the two security principles on which RAM-USB is built,
zero-trust and defense-in-depth.

\autoref{chap:requisiti} describes the methodology by which the requirements were gathered and traced,
and the use cases that derive from them.

\autoref{chap:architettura} presents the components that take part in registration and login.

\autoref{chap:autenticazione} follows the credentials along that path, from the public boundary to the
saved record, and describes the rules that each layer applies.

\autoref{chap:storage} follows the backup to its own isolated space on Storage-Service.

\autoref{chap:rete} describes how the mesh network enforces the reachability boundaries between the
components.

\autoref{chap:osservabilita} closes the lifecycle of a request with the
observability pipeline that makes it measurable.

\autoref{chap:testing} verifies the system thus built through the four levels of the
V-model, from unit testing to the acceptance checklist, and reports what running it on real hardware revealed.

\autoref{chap:deployment} describes how these components are actually deployed, across
Docker containers, Proxmox machines, and a public VPS.

\autoref{chap:sicurezza} rereads the entire system from the point of view of whoever attempts to compromise it,
verifying which guarantees hold even against its full compromise.

\autoref{chap:conclusioni} closes the thesis by weighing the results obtained against the limitations encountered,
and indicating the most relevant future work.
